\section{The main algorithm}

In order to tackle the minimum generalized conductance problem, our first objective is to further simplify the denominator $w\br{S, V}\cdot w\br{\bar{S}, V}$ into a linear function. We may without loss of generality assume that $w\br{S, V}\le w\br{\bar{S}, V}$.
Consequently, $w\br{V}/2\leq w\br{\bar{S}, V} \leq w\br{V}$ and thus, minimizing $\psi_w\br{S}$ is equivalent to minimization of
\[h_w\br{S}=\frac{c\br{S, \bar{S}}}{w\br{S, V}}, \quad \text{subject to} \quad w\br{S, V}\le w\br{\bar{S}, V}\] up to a constant factor loss. 
\DD{TODO: jeszcze moze dodac intuicje o tej stalej.}

However, for our purposes we will slightly relax the assumption that $w\br{S, V}\le w\br{\bar{S}, V}$ and in turn we will only require that $w\br{S, V}\le \C\cdot w\br{\bar{S}, V}$ for some appropriate constant $\C>1$.
\DD{$c$ juz jest zajeta na capacity, wiec zmienilem stala na $\C$ (uwaga: jest to makro, wiec mozna przedefiniowac).}
This still preserves the approximation ratio up to a factor dependent on $\C$ by a similar discussion.

As an internal tool used in this section we will utilize the generalized $k$-multicut problem, which is as follows.
Given a graph $G=(V,E,c,w)$ with capacities and demands, the \emph{generalized $k$-multicut problem} is to find a minimal-cost set of edges $C\subseteq E$, such that the demand cut by $C$ is at least $k$.

We now introduce a reduction\footnote{We use the word `reduction' since both $S_1$ and $S_2$ are obtained via an algorithmic reduction to another combinatorial problem, either generalized sparsest cut (with restrictions) or the generalized $k$-multicut.} scheme used in this section.
The scheme is applied to an input graph $G=(V,E,c_G,w_G)$ and it computes independently two candidate sets $S_1$ and $S_2$ via two different methos.
The returned solution is $S_i$ that minimizes $\psi_w$.
Formally, the scheme is the following:
\begin{enumerate}
  \item \label{it:S1} Compute a candidate set $S_1$:
  \begin{enumerate}[label=(\theenumi\alph*)]
    \item\label{it:S1:Gprime} Build an auxiliary graph $G'=(V',E',w_{G'},c_{G'})$ by adding a universal vertex $r$, i.e. $V'=V\cup\brc{r}$ and $E'=E\cup\brc{ru\colon u\in V}$, setting the capacities to $c_{G'}\br{r,u}=0$ for each $u\in V$ and $c_{G'}\brc{u,v}=c_G\brc{u,v}$ for each $u,v\in V$, and the demands to $w_{G'}\br{r,u}=\sum_{v\in V}w_G\br{u,v}$ for each $u\in V$ and $w_{G'}\br{u,v}=0$ for each pair $u,v\in V$.
    \item\label{it:S1:K} Set $K\gets \frac{4}{3}\cdot w_G\br{V}$.
    \item\label{it:S1:alpha} Run an $\alpha$-approximation algorithm for
    \[
      \min_{S\subseteq V'} \phi_w\br{S}=\frac{c_{G'}\br{S,\bar S}}{w_{G'}\br{S,\bar S}}
      \quad\text{subject to}\quad 0<w_{G'}\br{S,\bar S}\le K.
    \]
    and denote the returned cut by $S_1$. \DD{tutaj zamiast opisywac slownie problem odniosl bym sie, to frazy, ze rozwiazujemy generalized sparsest cut i podac dla jakiej instancji (tak, widze, ze bedzie extra ``subject to''. Jesli to ``subject to'' (nie wiem gdyż czytam liniowo) pojawialoby sie wiecej razy, to pewnie warto nazwac, a jak nie to niech zostanie tak, gdyz jest czytelnie; na razie zostawiam taki marker do oceny pozniej.}
  \end{enumerate}
  \item \label{it:S2} Compute a candidate set $S_2$:
  \begin{enumerate}[label=(\theenumi\alph*)]
    \item Run a $\beta$-approximation algorithm for the generalized $k$-multicut with $k=w_G\br{V}/3$ to obtain a multicut $C$.
    \item Build a complete graph $R$ on components of $G\setminus C$, with edge weights $w_R\br{H_1,H_2}=w_G\br{H_1,H_2}$ for any pair $H_1$ and $H_2$, where $H_1$ and $H_2$ are the vertex set of the components.
    \item Apply rounding: set $W=\frac{w\br{V\br{R}}}{4\cdot \spr{E\br{R}}}$ and replace the weight of each edge \DDfix{$e\in E\br{R}$} by $w'\br{e}=\fl{w\br{e}/W}$.
    \item Run a local search for max-cut on $(R,w')$: start with $(A,B)=(\emptyset,V\br{R})$, repeatedly move a single vertex from $B$ to $A$ if it improves the cut, and stop at a local optimum, i.e., when a single vertex transfer from $B$ to $A$ does not improve the cut.
    \item Map the resulting cut of $R$ back to a cut $S_2$ of $G$.
  \end{enumerate}
  \item Return the better of $S_1,S_2$, i.e., $\argmin_{X\in\brc{S_1,S_2}}\psi_w\br{X}$.
\end{enumerate}

\begin{theorem}
  The above algorithm is an $O\br{\alpha+\beta}$-approximation algorithm for the minimum generalized conductance problem.
\end{theorem}
\begin{proof}

Let $S^*$ be an optimal solution with $w\br{S^*, V}\le w\br{\bar{S}^*, V}$. Our analysis consists of two separate cases depending on the size of $w\br{S^*, \bar{S}^*}$.
We show that no matter what the actual optimal value is, at least one of the $S_i$ gives a good approximation ratio.
Note that the exact choice of constants in our construction is somewhat arbitrary and we do not attempt to optimize them.
\begin{lemma}
  Assume that $w\br{S^*, \bar{S}^*}\leq w\br{V}/3$. %Moreover, let $\alpha$ be an approximation ratio of some approximation algorithm for the problem of minimizing $\phi_w\br{S}$ in $G$ subject to $w_{G}\br{S, \bar{S}}\leq K$.
  Then, $S_1$ is an $O\br{\alpha}$-approximate solution to the minimum generalized conductance problem.
\end{lemma}
\begin{proof}
%To obtain a good approximation for $h\br{S}$, we proceed as follows. Firstly, we create an auxiliary graph $G'$, such that $V'=V\cup \brc{r}$. We set $c\br{r, u}=0$ for every $u\in V$ and we set $w\br{r, u}=\sum_{v\in V} w\br{u, v}$ for every $u\in V$.
Recall the auxiliary graph $G'$ defined in \ref{it:S1:Gprime}.
Take a cut $S\subseteq V$, assuming without loss of generality that $r\in \bar{S}$. We have that:
\[
w_{G}\br{S, V}\leq w_{G}\br{S, V}+w_G\br{S}= w_{G'}\br{S, \bar{S}}\leq 2\cdot w_{G}\br{S, V}
\]
\DD{
czy tu nie powinno byc?:
\[
w_{G}\br{S, V} = w_{G}\br{S, \bar{S}}+w_G\br{S}= w_{G'}\br{S, \bar{S}}\leq 2\cdot w_{G}\br{S, V}
\]
a jesli tak, to da nierownosc zachodzi trywialnie?
}
where the equality is due to the fact that for any $S\subseteq V\br{G'}$, the quanitity $w\br{S, \bar{S}}$ double counts the contribution of edges between the vertices of $S$ in $w_{G}\br{S, V}$.

\DD{tutaj na razie przerywam ten lemat, gdyz mam klopot z czytaniem: ten opis opiera sie sporo na intuicjach stojacych za budowa algorytmu zamiast po prostu policzyc to co jest. Na przyklad pierwszy paragraf budowal graph $G'$ --- dodalem jedno zdanie, zeby czytelnik sobie przypomnial punkt \ref{it:S1:Gprime}. Nastepny paragraf buduje intuicje, ale wydaje sie, ze mowi tylko cos w rodzaju ``Recall the definition of $K$ in \ref{it:S1:K}.''.}
It would be desireable to minimize the value of $\phi_w\br{S} = \frac{c\br{S, \bar{S}}}{w\br{S, \bar{S}}}
$ in this new graph. However, it is also easy to see, that in this new instance a trivial cut $S=V$ has cost $0$. Moreover, we demand $w_G\br{S, V}$ to be roughly smaller then $w_G\br{\bar{S}, V}$. To achieve both ends we thus add the constraint that $w_{G'}\br{S, \bar{S}}\leq K$ for $K=\frac{4}{3}\cdot w_G\br{V}$.

\DD{czemu odwolujemy sie do niezdefiniowanego $h(S)$ w dowodzie?}
We now want to argue that any $\alpha$-approximation of $\phi_w\br{S} $ subject to $w_{G'}\br{S, \bar{S}}\leq K$ in $G'$ is also a good approximation of $h\br{S}$ in $G$. Firstly, observe that for any optimal solution $S^*$ for $h\br{S}$, where $w\br{S^*, V}\leq w\br{\bar{S}^*, V}$ we have that
\[
w_G\br{S^*, V}+w_G\br{\bar{S}^*, V}=w_G\br{V}+w_G\br{S^*, \bar{S}^*}\leq \frac{4}{3}\cdot w_G\br{V}
\]
by our assumption that $w\br{S^*, \bar{S}^*}\leq w\br{V}/3$. As a consequence we get that $w_G\br{S^*, V}\leq \frac{2}{3}w\br{V}.$ In order to move this cut to $G'$ we append $r$ to $\bar{S}^*$.
This gives 
\[w_{G'}\br{S^*, \bar{S}^*}\leq 2\cdot w_G\br{S^*, V}\leq \frac{4}{3}\cdot w_G\br{V}=K.
\]
Therefore, we get that $\br{S^*, \bar{S}^*}$ is a valid solution for $\phi_{G'}\br{S}$ and as a consequence $\phi_{G'}\br{S}\leq h_{G}\br{S}$. Now we want to show an opposite direction, namely, that any solution to the problem of minimizing $\phi_w\br{S}$ in $G'$ subject to $w_{G'}\br{S, \bar{S}}\leq K$ gives a solution for $G$, such that $w_G\br{\bar{S}^*}\geq w\br{V}/6$. By the previous discussion, this yields an $O\br{\alpha}$-approximation for $h\br{S}$ and thus also for $\psi\br{S}$. Let $\br{S_1, \bar{S}_1}$ be the returned partition. 
There are 3 cases:
\begin{enumerate}
  \item $w_{G}\br{S_1, V}\leq \frac{5}{6}\cdot w_G\br{V}$. This means that  \[w_G\br{\bar{S}_1, V}\geq w_G\br{V}-w_{G}\br{S_1, V}\geq w_G\br{V}/6.
  \]
  \item $w_{G}\br{S_1, V}\geq \frac{5}{6}\cdot w_G\br{V}$. There are two subcases:
  \begin{itemize}
    \item $w_G\br{S_1}\geq w_G\br{V}/2$. We get that 
    \[w_G\br{S_1, V}=w_{G'}\br{S_1, \bar{S}_1}-w_G\br{S_1}\leq K-w_G\br{S_1} \leq \frac{4}{3}\cdot w_G\br{V}-w_G\br{V}/2 = \frac{5}{6}\cdot w_G\br{V}.
    \]
    This further means that 
    \[w_G\br{\bar{S}_1, V}\geq w_G\br{V}-w_G\br{S_1, V}\geq w_G\br{V}/6.
    \]
  \item $w_G\br{S_1, \bar{S}_1}\geq w_G\br{V}/3$. Here, we trivially get that $w_G\br{\bar{S}_1, V}\geq w_G\br{S_1, \bar{S}_1}\geq w_G\br{V}/3$.
  \end{itemize}
  It should be noted that at least one of the two above conditions must hold since $w_G\br{S_1}+w_G\br{S, \bar{S}_1}=w_{G}\br{S_1, V}$
\end{enumerate}
The lemma follows.
\end{proof}

Now we analyze, the second case. Here, we use an approximation ratio for  which is as follows: \DD{tutaj skleic...} This problem admits an $\beta=O\br{\log n}$-approximation algorithm by combining \cite{AUnifiedApproachToApproximatingPartialCoveringProblems} with \cite{OptimalHierarchicalDecompositionsForCongestionMinimizationInNetworks}.
\DD{Tw. mowi o tym, ze dostajemy $O(\alpha+\beta)$-approx i nie ma znaczenia ile one wynosza. Jesli chcemy tutaj rozkniniac, ze sa $\log n$-app. to by bylo zasadne, gdyby tw. to mowilo w tresci. Nie sugeruje tego, gdyz konstrukcja black-boxowa jest dobra: tw mowi o arbitralnym $\alpha+\beta$ i dopiero gdzies pozniej dywagacje jakie $\alpha$ i $\beta$ sa osiagalne i nastepnie finalny wspolczynnik. Natomiast poprzednie zdanie jest zbedne dla tego dowodu (moze to umnkelo, ale jestesmy w srodku dowodu twierdzenia).}
\begin{lemma}
  Assume that $w\br{S^*, \bar{S}^*}\geq w\br{V}/3$.
  %Moreover, let $\beta$ be an approximation ratio of some approximation algorithm for the (generalized) $k$-multicut problem.
  \DD{usunalem zalozenie o $\beta$, gdy to juz jest zdefiniowane w algo, a jesli chcemy w claimach przywolac, to jesli lematy sa zagniezdzone, to raczej w twierdzeniu.}
  Then, $S_2$ is an $O\br{\beta}$-approximate solution to the minimum generalized conductance problem.
\end{lemma}
\begin{proof}
  In this case, we have that $w\br{S^*, V}, w\br{\bar{S}^*, V}=\Theta\br{w\br{V}}$ \DD{w sumie tego nie widac, wiec dobrze byloby rozpisac/uzasadnic chocby jednym zdaniem} and therefore with a constant factor loss, it suffices to minimize $c\br{S, \bar{S}}$ subject to $w\br{S, \bar{S}}\geq w\br{V}/3$. \DD{tutaj dobrze byloby byc dokladniejszym co mamy na mysli: czy mamy const-factor loss wzgledem najlepszego rozwiazania, ktore spelnia $w\br{S, \bar{S}}\geq w\br{V}/3$ (i moze takie rozwiazanie jest gorsze od optimum bez tego ograniczenia), czy moze chodzi o to, ze zakladamy, ze input instance jest szczegolna i spelnia $w\br{S, \bar{S}}\geq w\br{V}/3$ i wtedy mamy const-factor loss wzgldem prawdziwego optimum dla takiej instancji?}
  Notably, this task is not equivalent to the $k$-multicut problem since $G\setminus C$ might consist of more then two components. To remedy this issue we proceeed as follows:
  \DD{mysle, ze raczej nalezaloby przerobic ten dowod, gdyz jest bardziej opisem algorytmu, natomiast klarowniej byloby tylko zrobic rachunki dajac ref-y do punktow algorytmu (zaproponowalem tam labele), aby czytelnik mogl spojrzec na algorytm i widziec ktory jego krok jest w danym momencie przeliczany.}
  We find an $\beta$-approximate multicut $C$. Then, we build a complete graph $R$, on vertex set consisting of components of $G\setminus C$. For any $H_1, H_2\in G\setminus C$, we set a weight $w_R\br{H_1, H_2}=w_G\br{H_1, H_2}$ (be weary of the fact that $H_1, H_2$ are vertices of $R$, but subsets of vertices of $G$). We now wish to solve a max-cut problem on $R$ in order to partition components of $G\setminus C$ in such a way to preserve a significant amount of demand being cut. Given a graph $R$, with weights $w$, the max-cut is to find a partition $\br{S_R, \bar{S_R}}$ of $V\br{R}$, such that $w\br{S_R, \bar{S_R}}$ is maximized. Although max-cut is NP-hard, it is well known that for any graph, there exists a cut $\br{A, B}$, such that $w\br{A, B}\geq w\br{V\br{R}}/2$. It is easy to see, that such a solution is $2$-approximation for max-cut, however for our purposes the crucial property is that since the demand cut by $C^*$ is at least $w_R\br{V\br{R}}=w_G\br{V}/3$, finding $(A, B)$ produces a cut of $G$ cutting at least $w\br{V}/6$ amount of demand which is sufficient for our purposes.
  
  To find an appropriate cut, we need an additional detail regarding the runtime of the our algorithm. There exists a local search policy which finds $\br{A, B}$ however its running time is pseudopolynomial in the weights.
  \DD{tutaj widac o co mi chodzilo: dowod musi uzupelniac algorytm. Parafrazujac o co mi chodzi: albo w typowym schemacie opis algorytmu i w dowodzie tylko analiza, albo zwyczajowo lemma mowi, ze cos da sie obliczyc i w dowodzie jest zaszyty algorytm. Tutaj mamy troche mix, oraz powtorzenia. Powtorzenia, gdyz sa zdania, ktore mowia co algorytm robi po kolei. Z drugiej strony sa uzupelnienia algorytmu. To drugie chyba jeszcze mogloby byc, choc pwnie lepiej byloby wyekstrachowac jakis lemat o local search mowiacy jak to zrobic i w jakim czasie i dac to jako cos niezaleznego poza tymi lematami i tym twierdzeniem (globalny lemat). Dodatkowe zamieszanie jest takie, ze analizujemy czas dzialania, wiec w tresci tw. trzebaby dodac, ze algorytm jest wielomianowy. Albo: ten globalny lemat powyzej moze moglby mowic, ze powyzszy schemat algorytmiczny dziala w wielomianoym czasie i tam przerzucic te kwestie zlozonosciowe.}
  To mitigate this inconvenience we perform a knapsack-like rounding step first. Let $W=\frac{w\br{V\br{R}}}{4\cdot \spr{E\br{R}}}$. For any $e\in E\br{R}$ we set $w'\br{e}=\fl{\frac{w\br{e}}{P}}$. Let $\br{A', B'}$ be the cut in this new instance for which $w'\br{A', B'}\geq w'\br{V\br{R}}/2$. For any edge $e$, $W\cdot w'\br{e}$ might be smaller then $w\br{e}$, but by no more then by $W$ (additively). Thus, by summing up over all edges cut by $\br{A, B}$ we get that:
\[
w\br{A, B}-W\cdot w'\br{e}\leq \spr{E\br{A, B}}\cdot W\leq \spr{E\br{R}}\cdot W.
\]
We have:
\[
w\br{A', B'}\geq W\cdot w'\br{A, B}\geq w\br{A, B} - \spr{E\br{R}}\cdot W = w\br{A, B}-w\br{V\br{R}}/4\geq w\br{V\br{R}}/4
\] 
We conclude that finding $\br{A', B'}$ gives us a cut of $R$ cutting at least $w\br{V\br{R}}/4$ amount of demand, which is still sufficient for out purposes.
Moreover, we have that 
\[
w'\br{V\br{R}}=\sum_{e\in E\br{R}}w'\br{e}\leq \frac{1}{W}\cdot \sum_{e\in E\br{R}}w\br{e}=4\spr{E\br{R}}^2
.\]
Hence, there are polynomially many values, that any cut of $\br{R, w'}$ can take, which will certify the polynomial runtime of our procedure.

In order to find $(A',B')$, we use the aforementioned local-search heurstic which is as follows: Start with $A=\emptyset, B=V\br{R}$. At any iteration, consider each vertex $v\in V\br{R}$, and if swapping the assignment of $v$ between $A$ and $B$ results in a better cut, perform the swap. Proceed until no such improvement is possible. For any vertex $v\in V\br{R}$, denote by $w'\br{v}$ its contribution to $w'\br{A', B'}$. We have the following:
\begin{align*}
2w'\br{A', B'}&=\sum_{v\in A'}w'\br{v}+\sum_{v\in B'}w'\br{v}
\\&\geq 
\sum_{v\in A'}\frac{w'\br{v, V\br{R}}}{2}+\sum_{v\in B'}\frac{w'\br{v, V\br{R}}}{2}
\\&=
\frac{1}{2}\sum_{v\in V\br{R}}w'\br{v, V\br{R}}=w'\br{V\br{R}}
\end{align*}
where the inequality is by the halting condition of our procedure.

The running time is polynomial, since every swap increases the weight of the cut, which as argued above can only take polynomially many different values.

Upon finding $\br{A', B'}$ we output the partition $\br{S_2, \bar{S}_2}$ of $V\br{G}$, such that vertices of members of $A'$ belong to $S_2$ and vertices of members of $B'$ belong to $\bar{S}_2$. We have that we cut at least $w\br{V}/12$ amount of demand, and the cost of the edges cut is $c\br{S_2, \bar{S}_2}\leq c\br{C}$, since $E\br{S_2, \bar{S}_2}\subseteq C$.
\end{proof}

By picking the best among the two created solution we ensure that the approximation ratio of the above procedure is upper bounded by $O\br{\alpha+\beta}$.
\end{proof}

Therefore it remains to obtain an approximation algorithm for the problem of minimizing $\phi_w\br{S}$ subject to $w\br{S, V}\leq K$, i.e., the algorithm used in step~\ref{it:S1:alpha}.
In what follows, we show how to do this using cut-sparsifier tree decompositions. We have the following proposition: \DD{czy to z literatury czy tuta? Jesli z lit. to trzebaby [??] dodać. Ogólnie dwuznaczna fraza jest ``subject to exact ...'': po pierwsze te claimy powinny mowic ``polynomial-time''; po drugie ten exact alg. for trees to trzebay napisac dla jakiego problemu, gdyz zaglujemy tutaj szeregiem problemow.}
\begin{proposition}
  There exists an $\alpha=O\br{\log n}$-approximation algorithm for the $\phi_w\br{S}$ subject to $w\br{S, V}\leq K$ problem
   for general graphs and an exact algorithm for trees.
\end{proposition}

It should be noted that the $O\br{\log n}$-approximation algorithm for the generalized $k$-multicut is also obtained using cut-sparsifiers \DD{domyslam sie, ze moze tutaj miec miejsce forward-referencja do innych miejsc. Trzeba dopilnowac, ze jesli uzywamy cos czego zrobimy to ze dla czytelnika to jest jasne i ze jasne jest to jakie brzmienie ma fakt, ktory uzywamy} in conjunction with an $O\br{1}$-approximation algorithm for trees. By combining all of the above discussions, we get the following corollary:
\begin{corollary}
  There is a polynomial-time $O(\log n)$-approximation algorithm for the problem of minimizing $\psi_w\br{G}$ on general graphs as well as an $O(1)$-approximation algorithm for the problem of minimizing $\psi_w\br{G}$ on trees.
\end{corollary}

A special kind of demand function is one, where for each vertex there exists a potential function $\pi\colon V\to \mathbb{N}$, such that for every $u,v\in V$, $w\br{u,v}=\pi\br{u}\cdot \pi\br{v}$. We call such a function a \emph{multiplicative function}. Notably, unit demand is multiplicative with $\pi\br{u}=1$ for every $u\in V$. For such demand functions, we have the following strengthening of approximation guarantees:
\begin{theorem}
  If the demand function $w$ is multiplicative, then there is a polynomial-time $O\br{\sqrt{\log n}}$-approximation algorithm for the problem of minimizing $\psi_w\br{S}$.
\end{theorem}
\begin{proof}
  Firstly, assume that $\pi\br{S}\leq \pi\br{\bar{S}}$. Since for multiplicative demands we have 
  \[
  w\br{A, B}=\sum_{u\in A}w\br{u}\cdot \sum_{v\in B}w\br{v}= \pi\br{A}\cdot \pi\br{B}
  \] 
  we get that $w\br{S, V}\leq w\br{\bar{S}, V}$ and moreover $w\br{S}\leq w\br{S, \bar{S}}$. Thus, minimizing $\psi_w\br{G}$ is up to a constant factor equivalent to minimizing 
  \[\frac{c\br{S, \bar{S}}}{w\br{S, \bar{S}}}=\frac{c\br{S, \bar{S}}}{\pi\br{S}\cdot \pi\br{\bar{S}}}.
  \]
  This can be then reduced to the min-ratio vertex cut problem which is as follows: Given a graph $G=\br{V,E,c,\omega}$, with a capacity vertex function $c\colon V\to \mathbb{N}$ and weight function $\omega\colon V\to \mathbb{N}$ find a partition $A, S, B$ of $V$ that minimizes 
  \[\frac{c\br{S}}{\omega\br{A\cup S}\cdot \omega\br{B\cup S}}
  \]
  such that there are no edges between $A$ and $B$. This problem admits an $O\br{\sqrt{\log n}}$-approximation algorithm \cite{ImprovedApproximationAlgorithmsForMinimumWeightVertexSeparators}. To reduce our instance to this problem we construct a new graph in which each edge $uv\in E$ is subdivided by a new vertex of capacity $c\br{u,v}$ and weight $0$ and each old vertex $u\in V$ is assigned capacity $0$ and weight $\pi\br{u}$. It is easy to see that the optimal solution for the min-ratio vertex cut problem on this new graph corresponds to the optimal solution for our problem on the original graph, and thus we get the desired approximation ratio. 
\end{proof}

